\chapter{Introducción}
\section{Motivación}
Un flujo de trabajo científico coordina tareas que consumen y producen archivos. Pegasus prepara el plan y HTCondor ejecuta sus unidades de trabajo, llamadas \textit{jobs}. El coordinador gestiona el flujo y recibe o envía sus archivos; W1 y W2 designan los dos nodos de ejecución (workers). En una cadena de tareas (\textit{pipeline}), cada etapa puede ejecutarse en un job separado o compartirlo con otras etapas. Esa decisión puede cambiar cuántas veces se entrega un archivo intermedio.

Esta tesis propone una forma de medir ese cambio y explicar qué entregas lo producen. Cuando el usuario agrupa dos etapas, puede dejar de reenviarse un intermedio al segundo job y conservarse su retorno al coordinador. Para distinguir esas operaciones hay que cotejar las dependencias, la ubicación de las tareas, el tamaño de los archivos y los registros de ejecución. En un caso pequeño se puede hacer la cuenta a mano. Aplicar el mismo criterio a varias configuraciones requiere un procedimiento explícito.

\section{Problema}
Un total de bytes carece de contexto si se desconoce qué archivos debían estar disponibles y en qué máquinas. Los registros de transferencia también pueden incluir ejecutables y otros archivos auxiliares. Para medir el volumen científico se necesita identificar los archivos que pertenecen al flujo y distinguirlos de los que requiere el sistema para ejecutarlo.

Agrupar tareas puede cambiar las entregas, y cambiar su ubicación también puede modificar las necesidades del flujo. La comparación debe separar esos dos efectos. Si un intermedio cruza entre workers, CondorIO exige su relevo por el coordinador. Si productor y consumidores permanecen en el mismo worker, sus contratos aún pueden declarar retornos y reenvíos que exceden esa cota.

\Needspace{5\baselineskip}
\textbf{Pregunta de investigación.} Dado un flujo de trabajo y la asignación real de sus tareas a máquinas, \textquestiondown{}qué cota inferior de movimiento de archivos imponen sus dependencias y el relevo obligatorio por el coordinador bajo CondorIO, y cómo se relaciona con las entregas completas corroboradas?

\section{Justificación}
La eficiencia del movimiento de datos (DME, por sus siglas en inglés) relaciona el volumen observado con las obligaciones que impone el flujo en su ubicación efectiva. Con las mismas necesidades, una reducción del observado indica que se entregaron menos bytes. Si cambian las necesidades por una nueva ubicación, la razón puede aumentar aunque el volumen entregado siga igual. Conservar las fuentes de ambos cálculos permite distinguir estas situaciones.

El aporte es una referencia por archivo y destino que permite interpretar el volumen entregado. La agrupación sirve para contrastarla: puede eliminar reenvíos, permitir compartir una entrada o conservar las mismas entregas. La evaluación distingue esos mecanismos y sus límites. Analyzer implementa el cálculo y vincula sus resultados con las fuentes que los sustentan.

\section{Objetivos}
\textbf{Objetivo general.} Diseñar, implementar y validar una herramienta que permita calcular una métrica para evaluar la eficiencia del movimiento de datos durante la ejecución de flujos de trabajo gestionados por Pegasus, utilizando los patrones process, pipeline, data aggregation, data distribution y data redistribution para evaluar la sensibilidad de la métrica bajo condiciones controladas de ejecución.

\textbf{Objetivos específicos.}
\begin{enumerate}
\item Caracterizar, a partir de la literatura y de la documentación técnica de Pegasus y HTCondor, las variables y fuentes de información necesarias para medir el movimiento de datos y la estructura de los patrones process, pipeline, data aggregation, data distribution y data redistribution.
\item Definir una métrica para evaluar la eficiencia del movimiento de datos, especificando sus variables, unidades, procedimiento de cálculo, supuestos, condiciones de aplicación y comportamiento esperado ante los patrones seleccionados.
\item Diseñar e implementar una herramienta que automatice la extracción, integración y procesamiento de la información requerida para calcular la métrica definida a partir de ejecuciones gestionadas por Pegasus.
\item Validar la métrica y la herramienta mediante ejecuciones controladas de los cinco patrones de movimiento de datos, documentando la corrección y consistencia del cálculo, la sensibilidad de la métrica, las condiciones de aplicación, el sobrecosto de instrumentación, las fortalezas y las limitaciones de la solución.
\end{enumerate}


\Needspace{5\baselineskip}
\section{Alcance}
La tesis abarca la especificación de DME, su cálculo con \Analyzer{} y su evaluación en flujos gestionados por Pegasus y HTCondor con CondorIO. El modelo distingue la cota que exigen el grafo y el relevo por el coordinador con una asignación dada del volumen de las entregas corroboradas. La herramienta integra el plan, el historial y los registros de transferencia, presenta el desglose por archivo y destino, y señala cuándo la evidencia es insuficiente.

La medición comprende las entradas científicas, los intermedios y los resultados declarados, desde las entradas disponibles en el coordinador hasta la recuperación de los resultados. La unidad de ubicación es la máquina virtual y los archivos deben conservar su identidad y tamaño. Se contabilizan entregas completas. Los registros de archivos auxiliares se usan para conciliar los contadores de transferencia. Los volúmenes se expresan en bytes y DME es su cociente adimensional. Para aplicar el procedimiento a otro flujo se requiere evidencia de sus archivos, ubicaciones y entregas bajo el mismo perfil de medición.

El bloque A reúne 24 condiciones y 72 ejecuciones elegibles. La validación histórica detallada analiza 45 ejecuciones de 15 condiciones en los cinco patrones. En un contraste separado se comparan L01, GL2 y GL3 con ubicación libre (n=3) y se documenta un piloto GL2; las asignaciones efectivas se recuperan del historial. Se conservan también cuatro ejecuciones históricas de Montage con ubicación W1/W2. Una campaña prospectiva separada completa 63 ejecuciones principales: seis por condición sintética y tres por variante de Montage con ambas tareas en W1. Doce pilotos válidos y un intento técnico fallido se documentan aparte. Se comparan agrupaciones y una intervención adicional en el retorno de intermedios, conservando los finales científicos. La duración es una respuesta secundaria. El costo de generar registros durante la ejecución queda por aislar.

\section{Contenido y relación con los objetivos}
El capítulo~\ref{ch:marco} presenta los patrones, conceptos y fuentes del primer objetivo. El capítulo~\ref{ch:diseno} define la métrica y sus condiciones de aplicación, correspondientes al segundo. El capítulo~\ref{ch:implementacion} documenta la herramienta del tercero. El capítulo~\ref{ch:validacion} aborda el cuarto mediante la comprobación del cálculo, la consistencia, el contraste libre y la campaña prospectiva de agrupación. Los capítulos finales presentan las conclusiones, el trabajo futuro y los materiales de reproducción.

\begin{table}[htbp]\centering\small
\caption{Evidencia asociada a los objetivos específicos}
\label{tab:objetivos_evidencia}
\begin{tabularx}{\linewidth}{@{}P{2.5cm}YY@{}}
\toprule Objetivo & Evidencia en la tesis & Qué se puede comprobar\\\midrule
OE1. Caracterización & Definiciones, cinco grafos y fuentes en el marco teórico y la implementación. & Qué cambia entre tareas, jobs, ubicaciones y archivos, y de dónde se obtiene cada variable.\\
OE2. Definición & Cota por archivo y destino con relevo, unidades, cobertura y estados. & Cuándo se puede calcular DME y cuándo el observado o la referencia explican el cambio.\\
OE3. Implementación & Arquitectura, contratos, diagnósticos y materiales de reproducción. & Cómo se obtienen desgloses verificables. La regla de relevo está implementada.\\
OE4. Validación & Cotejo histórico de 45 ejecuciones, nueve libres y un piloto; 63 principales prospectivas en doce condiciones, con pilotos separados y control de finales. & Correctitud, consistencia y sensibilidad bajo el perfil ensayado. Queda por aislar el costo de captura nativa.\\
\bottomrule
\end{tabularx}
\end{table}
